% =====================================================================
% Section 1 — Introduction
% =====================================================================

\section{Introduction}\label{sec:intro}

Spark-ignition (SI) engines remain technologically relevant in light
automotive applications, particularly in markets where cost pressure,
mechanical robustness, an established refuelling network, and a gradual
transition to lower-impact energy solutions coexist
\citep{heywood1988,stone2012}. Within this landscape, the search for
higher specific output and thermal efficiency does not depend solely on
electronic calibration or fuel selection: the geometric quality of the
intake system, of the combustion chamber, and of the valve train, together
with thermal management of the head, governs the practical envelope of
operation.

In high-performance configurations, where elevated rotational speed,
increased thermal loading, and stricter operational stability are required
simultaneously, two risks acquire a critical role: abnormal combustion
(knock) and valve-train failure. Both phenomena compromise reliability,
shrink the operational safety margin, and may evolve from a local
component issue into catastrophic damage to the engine. The present work
focuses exclusively on the second of these two risks. The companion
problem of knock at high compression ratio is addressed elsewhere and is
out of scope here.

% --- 1.1 The case for studying the valve train of the AP 1.8 -----------
The Volkswagen AP 1.8 platform was adopted as the reference engine for
two reasons. First, it occupies an established position in the Brazilian
automotive market, with a wide installed base and broad availability of
geometric data \citep[][\S~1.6]{porto2026dissertation}. Second, in the
context of a master's research project investigating a high-performance
cylinder head for the same platform, a specific operating regime of
10\,000~rpm was adopted as the design target. At that rotational speed
the inertial demand placed on the valve--cam contact pair is no longer a
local component concern: the dynamic stability of the moving valve mass,
the residual coil-bind clearance, the natural frequency of the spring
relative to cam harmonics, and the fatigue resistance of the spring all
become coupled, and a single insufficient margin can propagate into a
valve--piston collision.

% --- 1.2 Why FMEA/RPN alone is insufficient ----------------------------
The conventional Failure Mode and Effects Analysis (FMEA) approach
addresses such situations by ranking failure modes through the
multiplicative Risk Priority Number (RPN). This approach has known
limitations: subjectivity in the assignment of severity, occurrence, and
detectability scales; loss of resolution when multiple modes share a
similar RPN; and limited visibility of the causal chain that links
threats to consequences via barriers. The BowTie methodology addresses
these limitations by organising the analysis around a top event, with
threats and preventive barriers on the left side, and consequences and
mitigative barriers on the right side. This structure makes the role of
each barrier explicit and supports communication between design,
operation, and reliability domains.

% --- 1.3 Proposal of this paper ----------------------------------------
This paper proposes a quantitative-coupled BowTie analysis of the valve
train of a high-performance cylinder head operating at 10\,000~rpm in the
VW AP 1.8 engine. The qualitative bow diagram is enriched with a
quantitative layer derived from the analytical sizing of the intake
valve and of the valve springs. Each preventive barrier is associated
with a quantitative pass criterion (natural-frequency margin, coil-bind
margin, fatigue factor of safety) so that the BowTie diagram becomes both
a communication artefact and a design check.

% --- 1.4 Contributions -------------------------------------------------
The contributions of this paper are summarised as follows:
\begin{enumerate}
\item A BowTie diagram of the valve train of the AP 1.8 cylinder head at
10\,000~rpm, with the top event ``loss of follower contact between valve
and cam'' and the associated set of threats, preventive barriers,
consequences, and mitigative barriers.
\item A quantitative coupling layer that translates each preventive
barrier into a measurable indicator with a literature-anchored pass
criterion, derived from the analytical sizing of the valve and the
spring.
\item Identification of the critical preventive barrier for the operating
envelope adopted, with a comparison against a conventional FMEA/RPN
ranking of the same failure modes.
\item A reproducible methodology that can be transposed to other
high-performance cylinder heads operating at high rotational speed.
\end{enumerate}

% --- 1.5 Positioning of the contribution -------------------------------
\noindent\textbf{Positioning of the contribution.} This work is a
\emph{prospective, design-stage} reliability assessment of a
high-rpm valve train; it is not a forensic post-failure analysis.
The manuscript identifies the relevant failure mechanisms of the
high-rpm valve train (valve float, bounce, approach to coil bind,
surge resonance, fatigue) and proposes barrier-centric preventive
actions, organising the analytical sizing of the valve and the
springs around the BowTie methodology. The novelty in relation to
the international literature is domain-specific and incremental, not
methodological in the broad sense: quantitative extensions of the
BowTie diagram have been previously proposed
(\citealp{khakzad2013bowtiebayesian}, via Bayesian-network mapping;
LOPA, fault-tree-coupled BowTie, and DyPASI in process safety),
and the present work does not claim to introduce a new general
method. The original contribution is the explicit mapping of each
preventive barrier to a deterministic indicator derived from the
engineering sizing of the intake valve and the dual valve springs of
the AP 1.8 head, with a literature-anchored pass criterion for each,
and the identification of the resulting critical barriers under the
operating envelope adopted. The procedure is qualitatively validated
against four documented valve-train failure cases from the recent
literature
\citep{velezgodino2019overheadvalvetrain,soffritti2018wornvalvetrain,xing2021valvespringEAC,zhao2023camshaftV6},
showing that each reported root cause maps to a specific preventive
barrier of the BowTie diagram and would have been flagged at the
design stage by the corresponding indicator. An empirical
verification on a dedicated test rig is identified as the priority
direction for future work, and a concrete validation roadmap (see
Section~\ref{sec:conclusion}) specifies the test-rig instrumentation
and acceptance criteria. The claim that no published study applies
the BowTie methodology to the valve train of a high-rpm
spark-ignition cylinder head is supported by a search of the Scopus
and Web of Science databases conducted in April 2026 with the query
strings
\texttt{("BowTie" OR "bow-tie") AND ("valve train" OR "valve-train" OR "valve spring") AND ("spark ignition" OR "SI engine" OR "high rpm")},
which returned no relevant hits.

% --- 1.6 Paper outline -------------------------------------------------
Section~\ref{sec:background} reviews the relevant failure modes of the
valve train at high rpm and the BowTie methodology.
Section~\ref{sec:methodology} describes the system under study, the
operational premises, the qualitative BowTie construction, and the
quantitative-coupling procedure. Section~\ref{sec:results} reports the
populated BowTie diagram, the values of the quantitative indicators
for the AP 1.8 case, and the qualitative validation of the procedure
against four published valve-train failure cases.
Section~\ref{sec:discussion} discusses the implications for
cylinder-head design, compares the proposed approach with FMEA/RPN,
and states the limitations of the present analytical study.
Section~\ref{sec:conclusion} closes the paper and specifies the
experimental validation roadmap. The work complements the
post-failure case-study tradition of
\citet{velezgodino2019overheadvalvetrain,soffritti2018wornvalvetrain,xing2021valvespringEAC,zhao2023camshaftV6},
which informed the present formulation, by translating their lessons
into a structured prospective barrier-evaluation procedure applicable
at the design stage.
